\section{Conclusión}\label{sec:conclusion}

Desde una conexión domiciliaria de WOM en Quilicura, medí las rutas hacia Valparaíso, Cuba y Rusia con Scamper y con RIPE Atlas, en tres rondas y en los dos sentidos, 153 ejecuciones en total 

\textbf{¿Qué método descubre más interfaces?} No hay un único ganador, depende del destino y del criterio, y conviene mirar siempre el costo en sondas. La tabla resume el mejor método de Scamper, el mejor protocolo de RIPE y lo que agregó tracelb.


\begingroup
\fontsize{9}{11}\selectfont
\addtocounter{table}{-1}
\begin{longtable}{p{2.2cm}p{4.7cm}p{4.6cm}p{3.3cm}}
\toprule
\textbf{Destino} & \textbf{Scamper} & \textbf{RIPE Atlas} & \textbf{tracelb} \\ \midrule
\endfirsthead
\toprule
\textbf{Destino} & \textbf{Scamper} & \textbf{RIPE Atlas} & \textbf{tracelb} \\ \midrule
\endhead
Valparaíso & icmp, icmp-paris, udp y udp-paris: 9 IP, 9 a 10 sondas & ICMP, TCP y UDP: 9 IP, 27 sondas & sin IP adicionales \\
La Habana & icmp e icmp-paris: 13 IP, estab. 0.949, 22 sondas & TCP: 14 IP, estab. 0.730, 57 sondas & 3 IP adicionales con udp-dport (214 sondas) \\
Rusia & icmp e icmp-paris: 12 IP, estab. 0.944, 18 a 19 sondas & TCP: 11 IP, estab. 0.840, 51 sondas & 1 IP adicional con udp-dport (204 sondas) \\
\bottomrule
\end{longtable}
\endgroup

En Valparaíso, con una ruta de nueve saltos, todos los métodos que reciben respuesta empatan. En Cuba (Habana, 14 contra 13) y mejor cobertura en los dos destinos lejanos (82,4\%), pero con el doble o el triple de sondas y menos estabilidad entre rondas. En Rusia no superó a Scamper en IP únicas (11 contra 12). tracelb mostró interfaces adicionales solo con udp-dport, solo en La Halana y en Rusia, a un costo cercano a diez veces el de un trace y sin repetirse en todas las rondas. Para describir una ruta, ICMP con Scamper ofrece la mejor relación entre información y costo, mientras que tracelb y los varios paquetes por salto de RIPE sirven sobre todo para detectar balanceo de carga. También hay que recordar que en Scamper los métodos TCP no registraron ninguna IP, pero eso se explica por la recepción y no por la red. \\

Al invertir el sentido cambia mucho. En los tres destinos casi no hay direcciones en común entre ida y vuelta (Jaccard entre 0.00 y  0.08), y en La Habana y en Rusia además son distintos los sistemas autónomos visibles, la ida sale por AS174 y la vuelta llega por AS3356 y AS3549, con las mismas interfaces finales para ambos destinos. Eso lo convierte en una evidencia de asimetría más fuerte que la mera diferencia de direcciones, que se esperaría incluso con un camino simétrico, porque cada sentido muestra la interfaz de entrada. En Valparaíso, en cambio, los segmentos de red se parecen y la diferencia podría deberse solo a qué interfaz se ve en cada sentido, aunque los huecos impiden afirmar simetría completa. Además, solo ICMP logró que la vuelta llegara hasta mi propia dirección. \\

Para reconocer las sondas en una captura, RIPE Atlas es la más fácil de identificar, sus paquetes UDP, ICMP y TCP llevan en el payload el texto \verb|http://atlas.ripe.net Atlas says Hi!|, en 48 bytes. Scamper no incluye texto, usa payloads de 16 bytes en cero con los dos primeros bytes variables en las variantes Paris y en tracelb, o no usa payload en TCP, y se reconoce por los campos que cambian (el puerto de destino que parte en 33435, el icmp.seq que parte en 0, el TTL que parte en 1 o el IPID que parte en 0x001 en tracelb).
En los casos UDP e ICMP, la suma de payload, checksum y campos de flujo se mantuvo constante, que es lo esperable con una compensación de checksum, pero esa regularidad también la produce cualquier sonda que cambie solo esos campos y no prueba por sí sola la intención del emisor. \\

Para la respuesta selectiva y TCP, la IP del destino ruso respondió a icmp y no a udp en dos repeticiones, y 200.0.16.25m en la ruta a Cuba, respondió a TCP en RIPE y no a ICMP, también en dos repeticiones. Los contrastes son repetibles, pero no identifican una regla de filtrado concreta, otra ruta, un balanceo o la pérdida de paquetes también los explicarían. Por otro lado, los ceros de Scamper con TCP no significan que la red bloquee TCP: en 36 ejecuciones aparecieron en la captura respuestas ICMP que Scamper no registró como saltos, de modo que esos ceros hablan de lo que la herramienta aceptó y no de lo que la red dejó pasar. \\

El principal límite es que todo lo anterior vale para una sola conexión, una ventana de tiempo acotada y tres destinos. Las IP son interfaces, no probé que interfaces pertenecen al mismo router, así que no estoy contando routers. Los ASN salen de asociaciones IP-prefijo y no son una ruta BGP medida. Los dos primeros destinos vienen de una campaña anterior a la de Rusia y no forman una instantánea simultánea. Un traceroute no mide velocidad de acceso ni demuestra censura, las diferencias entre Chile, Cuba y Rusia describen las rutas y los horarios medidos, no una clasificación de velocidades.\\

En resumen, lo principal es que la calidad de un Path Discovery depende tanto del método como de la captura, de la recepción local y de la trazabilidad. El caso de TCP lo muestra bien, la salida de la herramienta por sí sola, habría sugerido que esas rutas no existían, y solo al contrastarla con los paquetes se vio que las respuestas sí estaban. Para repetir el estudio conviene preservar los originales, contrastar los parámetros efectivos, verificar la continuidad de la sesión y comparar pares próximos en el tiempo. Como trabajo futuro, aplicaría resolución de alias para contar routers y no solo interfaces, repetiría las mediciones en varios horarios para separar la variación temporal de la estructural, y revisaría los WARTS de tracelb y de TCP para explicar los conteos de sondas idénticos entre destinos y los ceros de recepción.






